Este trabalho analisa o web-flickr, grafo da coleção SNAP em que cada vértice é uma imagem do Flickr e cada aresta liga duas imagens que compartilham metadados. A lista de arestas foi convertida em um grafo simples não direcionado, e sobre ele e a sua maior componente conexa calcularam-se, de forma exata, as medidas estruturais pedidas, que reproduzem as estatísticas declaradas pelo SNAP. Os algoritmos de busca em largura e em profundidade, verificação de Eulerianidade, Tarjan, Dijkstra, Bellman-Ford, Floyd-Warshall e Kruskal foram implementados, validados contra o networkx e cronometrados com intervalos de confiança, e os tempos observados foram comparados à complexidade teórica. O grafo apresenta indícios de small-world, com clusterização muito maior que a do grafo aleatório equivalente e comprimento médio dos caminhos próximo do dele. A lei de potência descreve só o extremo da cauda da distribuição de graus. A maior componente resiste à remoção aleatória de 5\% dos vértices e à remoção dos 5\% de maior grau, embora o ataque direcionado isole muito mais vértices da periferia. A descoberta mais interessante é que os vértices de maior grau ficam em grupos densos e redundantes entre si, e não são os elos que mantêm a rede conectada.
